\documentclass[11pt]{article}

\usepackage[margin=1in]{geometry}
\usepackage{amsmath,amssymb,amsthm}
\usepackage{array}
\usepackage[hidelinks]{hyperref}

\newtheorem{theorem}{Theorem}[section]
\newtheorem{lemma}[theorem]{Lemma}
\newtheorem{proposition}[theorem]{Proposition}
\newtheorem{corollary}[theorem]{Corollary}
\theoremstyle{definition}
\newtheorem{definition}[theorem]{Definition}
\newtheorem{remark}[theorem]{Remark}

\newcommand{\Cl}{\mathrm{Cl}}
\newcommand{\Spin}{\mathrm{Spin}}
\newcommand{\SO}{\mathrm{SO}}
\newcommand{\spinlie}{\mathfrak{spin}}
\newcommand{\solie}{\mathfrak{so}}
\newcommand{\Spec}{\mathrm{Spec}}
\newcommand{\ad}{\mathrm{ad}}
\newcommand{\tr}{\mathrm{tr}}
\newcommand{\Euc}{\mathrm{Euc}}
\newcommand{\rev}[1]{\widetilde{#1}}
\newcommand{\grade}[2]{\langle #1 \rangle_{#2}}

\title{\textbf{Disproof and Complete Geometric-Algebraic Characterization\\
of Rotor Logarithm-Sum Angle Additivity}\\
\large Resolution of Conjecture 00000007607}
\author{Formalized in Core Lean 4.19.0 with Exact Computational Verification}
\date{}

\begin{document}
\maketitle

\begin{abstract}
We resolve Conjecture \texttt{00000007607}, which asserts under ``the exponential law of
rotor formulas'' that for rotor generators in Euclidean geometric algebra: (I) the logarithm
sum of generators gives the composite rotation angle; (II) the spectrum of the generator sum
exhibits angle additivity; and (III) angle additivity is verified by the Lie bracket identity.
We prove that the conjecture is \textbf{false in general} across every Euclidean dimension
$n \ge 3$, while \textbf{true in dimension $n = 2$} (and more generally if and only if the
participating simple rotors share the same oriented rotation plane). Using the canonical
Clifford grade decomposition $B_1 B_2 = \grade{B_1 B_2}{0} + \frac{1}{2}[B_1,B_2] + B_1 \wedge B_2$
for bivectors $B_1, B_2 \in \spinlie(n) = \bigwedge^2 \mathbb{R}^n$, we establish exact
polynomial sum-of-squares identities governing rotor composition, generator-sum spectra, and
commutator norms over arbitrary ordered commutative rings. In $\Spin(3) \subset \Cl_{3,0}$,
composing two orthogonal $90^\circ$ ($\pi/2$) rotors produces a $120^\circ$ ($2\pi/3$) cyclic
rotation, whereas the exponential of the logarithm sum has rotation angle $\pi/\sqrt{2} \approx
127.28^\circ$ and naive angle additivity gives $180^\circ$ ($\pi$), separating all three
quantities simultaneously while $[e_{12}, e_{23}] = 2e_{13} \neq 0$ acts as the
Baker--Campbell--Hausdorff obstruction rather than a universal verification identity. Moreover,
in $\Spin(4) \subset \Cl_{4,0}^+$ we prove that even when the Lie bracket vanishes
($[e_{12}, e_{34}] = 0$) and $\exp(B_1)\exp(B_2) = \exp(B_1 + B_2)$ holds, the non-vanishing
exterior wedge $B_1 \wedge B_2 \neq 0$ produces a non-simple double rotation whose vector
spectrum is $\{\pm i\theta_1, \pm i\theta_2\}$ rather than single-angle additivity. All
algebraic characterizations, positive planar theorems, and counterexamples across all $n \ge 3$
are formally verified from first principles in Core Lean~4.19.0 with zero external dependencies.
\end{abstract}

\section{Introduction and Problem Formulation}

Conjecture \texttt{00000007607} is stated in \texttt{SOURCE.md} in parallel English and
Chinese formulations as follows:
\begin{quote}
\textbf{Definition:} The exponential law of rotor formulas. \textbf{Conjecture:} The
logarithm sum of generators gives the rotation angle; the spectrum of the sum is angle
additivity; and the verification of the angle additivity is the Lie bracket identity.
(\emph{rotor logarithm-sum angle additivity})
\end{quote}
The parallel Chinese statement similarly defines the exponential law of rotor formulas
(\emph{zhuanzi gongshi de zhishulv}) and conjectures that the logarithm sum of generators
yields the rotation angle (\emph{shengchengyuan de duishu he geichu xuanzhuanjiao}), the
spectrum of the sum is angle additivity (\emph{he de pu wei jiao jiaxing}), and the
verification of angle additivity is the Lie bracket identity (\emph{jiao jiaxing de yanzheng
wei Li kuohao hengdengshi}).

\subsection{Geometric Algebra and Spin-Rotor Preliminaries}

Let $\Cl_{n,0}(\mathbb{R})$ denote the Euclidean Clifford (geometric) algebra of $\mathbb{R}^n$
generated by an orthonormal basis $e_1, \dots, e_n$ subject to
\begin{equation}
  e_i^2 = 1 \quad (1 \le i \le n), \qquad
  e_i e_j = -e_j e_i \quad (1 \le i \neq j \le n).
\end{equation}
As a vector space, $\Cl_{n,0}(\mathbb{R}) = \bigoplus_{k=0}^n \bigwedge^k \mathbb{R}^n$ has
dimension $2^n$, with grade-$k$ projection denoted $\grade{X}{k}$ and basis blades
$e_{i_1 \cdots i_k} = e_{i_1} \cdots e_{i_k}$ ($1 \le i_1 < \cdots < i_k \le n$).
\emph{Reversion} $X \mapsto \rev{X}$ is the unique anti-automorphism ($\rev{XY} = \rev{Y}\rev{X}$)
fixing vectors ($\rev{e_i} = e_i$); on grade $k$ it acts by $(-1)^{k(k-1)/2}$, fixing grades
$0, 1, 4$ and negating grades $2, 3$.

On the bivector space $\bigwedge^2 \mathbb{R}^n$, the positive-definite Euclidean inner product
for which $\{e_{ij} : 1 \le i < j \le n\}$ is orthonormal is given by
\begin{equation}\label{eq:bivec-inner}
  \langle X, Y \rangle_{\Euc} = \grade{X \rev{Y}}{0} = -\grade{X Y}{0}, \qquad
  \|X\|_{\Euc}^2 = \langle X, X \rangle_{\Euc} = -\grade{X^2}{0}.
\end{equation}
Equipped with the commutator Lie bracket $[X, Y] = X Y - Y X$, the bivector space forms the
spin Lie algebra $\spinlie(n) = \bigwedge^2 \mathbb{R}^n$. Its exponential image inside the
even subalgebra $\Cl_{n,0}^+(\mathbb{R})$ is the spin group
\begin{equation}
  \Spin(n) = \{ R \in \Cl_{n,0}^+(\mathbb{R}) : R \rev{R} = 1, \;
  R v \rev{R} \in \mathbb{R}^n \text{ for all } v \in \mathbb{R}^n \},
\end{equation}
whose elements are called \emph{rotors}. The double-covering group homomorphism
$\rho : \Spin(n) \twoheadrightarrow \SO(n)$, $\rho(R)(v) = R v \rev{R}$, differentiates at
$R = 1$ to the Lie algebra isomorphism $\ad : \spinlie(n) \xrightarrow{\sim} \solie(n)$,
$\ad_B(v) = [B, v] = B v - v B$.

A \emph{unit simple bivector} $u \in \bigwedge^2 \mathbb{R}^n$ represents an oriented $2$-plane
($u = a \wedge b = a b$ for orthonormal vectors $a, b \in \mathbb{R}^n$) and satisfies
$u^2 = -1$ and $u \wedge u = 0$. By Euler's power-series expansion of
$\exp(X) = \sum_{m=0}^\infty X^m / m!$, a \emph{simple rotor} rotating the plane of $u$ by
angle $\theta \in \mathbb{R}$ is
\begin{equation}\label{eq:euler-rotor}
  R(\theta, u) = \exp\!\left(-\frac{\theta}{2} u\right)
  = \cos\frac{\theta}{2} - \sin\frac{\theta}{2}\, u,
  \qquad B = -\frac{\theta}{2} u \in \spinlie(n).
\end{equation}
Here $B$ is the logarithmic bivector generator of $R(\theta, u)$ (its principal logarithm for
$|\theta| < 2\pi$), with $\|B\|_{\Euc} = |\theta|/2$. In the complexified spinor representation
$L_B(X) = B X$, $L_B^2 = -\frac{\theta^2}{4} I$ has spectrum
$\Spec(L_B) = \{\pm i\theta/2\}$; in the vector representation on $\mathbb{C}^n$,
$\ad_B \in \solie(n)$ rotates the plane of $u$ by $\theta$ and annihilates its orthogonal
complement, giving $\Spec(\ad_B) = \{\pm i\theta\} \cup \{0\}^{(n-2)}$.

\subsection{The Three Clauses of Conjecture 00000007607}

Given two simple rotors $R_1 = \exp(B_1)$ and $R_2 = \exp(B_2)$ with logarithmic generators
$B_k = -\frac{\theta_k}{2} u_k$ ($\theta_k \in (0, \pi]$, $u_k^2 = -1$), Conjecture
\texttt{00000007607} asserts three linked properties:
\begin{enumerate}
  \item[\textbf{(I)}] \textbf{Logarithm-sum rotation angle:} Under the exponential law
    $\exp(B_1)\exp(B_2) = \exp(B_1 + B_2)$, the logarithm sum $B_1 + B_2$ generates the
    composite rotor $R_1 R_2$ and determines its composite rotation angle $\Theta$, yielding
    $\Theta = 2\|B_1 + B_2\|_{\Euc} = \theta_1 + \theta_2$.
  \item[\textbf{(II)}] \textbf{Spectral angle additivity:} The spectrum of the generator sum
    $B_1 + B_2$ (in either the spinor representation $L_{B_1 + B_2}$ or the vector
    representation $\ad_{B_1 + B_2}$) equals that of a single rotation of angle
    $\theta_1 + \theta_2$, i.e.\ $\|B_1 + B_2\|_{\Euc} = \frac{\theta_1 + \theta_2}{2}$ and
    $\Spec(\ad_{B_1 + B_2}) = \{\pm i(\theta_1 + \theta_2)\} \cup \{0\}^{(n-2)}$.
  \item[\textbf{(III)}] \textbf{Verification via Lie bracket identity:} Angle additivity is
    verified by a universal Lie bracket identity on rotor generators in $\spinlie(n)$.
\end{enumerate}

We show below that \textbf{all three clauses hold identically when $n = 2$} (Theorem~\ref{thm:n2}),
and \textbf{all three clauses fail in every dimension $n \ge 3$}
(Theorems~\ref{thm:spin3-iff}--\ref{thm:universal-n}).

\section{Canonical Clifford Bivector Grade Decomposition}

The complete behavior of rotor composition across all dimensions is governed by a single
structural decomposition of the geometric product of two bivectors.

\begin{lemma}[Bivector Grade Decomposition]\label{lem:grade-decomp}
For all bivectors $B_1, B_2 \in \bigwedge^2 \mathbb{R}^n$ in $\Cl_{n,0}(\mathbb{R})$,
\begin{equation}\label{eq:grade-decomp}
  B_1 B_2 = -\langle B_1, B_2 \rangle_{\Euc}
  + \frac{1}{2}[B_1, B_2] + B_1 \wedge B_2,
\end{equation}
where $-\langle B_1, B_2 \rangle_{\Euc} = \grade{B_1 B_2}{0} \in \mathbb{R}$,
$\frac{1}{2}[B_1, B_2] = \grade{B_1 B_2}{2} \in \bigwedge^2 \mathbb{R}^n$, and
$B_1 \wedge B_2 = \grade{B_1 B_2}{4} \in \bigwedge^4 \mathbb{R}^n$. Consequently,
\begin{equation}\label{eq:sum-sq-general}
  (B_1 + B_2)^2 = B_1^2 + B_2^2 - 2\langle B_1, B_2 \rangle_{\Euc} + 2(B_1 \wedge B_2).
\end{equation}
\end{lemma}

\begin{proof}
Multiplying two basis $2$-blades $e_{ij}$ and $e_{kl}$ produces grade $0$ if $\{i,j\}=\{k,l\}$,
grade $2$ if $|\{i,j\} \cap \{k,l\}| = 1$, and grade $4$ if $\{i,j\} \cap \{k,l\} = \emptyset$.
Thus $B_1 B_2 = \grade{B_1 B_2}{0} + \grade{B_1 B_2}{2} + \grade{B_1 B_2}{4}$. Since reversion
negates grade $2$ and fixes grades $0$ and $4$, we have
$B_2 B_1 = (-B_2)(-B_1) = \rev{B_2}\rev{B_1} = \rev{B_1 B_2}
= \grade{B_1 B_2}{0} - \grade{B_1 B_2}{2} + \grade{B_1 B_2}{4}$.
Taking half-sum and half-difference yields $\frac{1}{2}(B_1 B_2 - B_2 B_1) = \grade{B_1 B_2}{2}$
and $\frac{1}{2}(B_1 B_2 + B_2 B_1) = \grade{B_1 B_2}{0} + \grade{B_1 B_2}{4}
= -\langle B_1, B_2 \rangle_{\Euc} + B_1 \wedge B_2$, proving \eqref{eq:grade-decomp} and
\eqref{eq:sum-sq-general}.
\end{proof}

Each grade piece in \eqref{eq:grade-decomp} controls a distinct aspect of Conjecture
\texttt{00000007607}:
\begin{itemize}
  \item \textbf{Grade 2 ($\frac{1}{2}[B_1, B_2]$):} Controls commutativity and the exponential
    law $\exp(B_1)\exp(B_2) = \exp(B_1 + B_2)$ via the Baker--Campbell--Hausdorff formula.
  \item \textbf{Grade 4 ($B_1 \wedge B_2$):} Controls simplicity of the sum $(B_1 + B_2) \wedge
    (B_1 + B_2) = 2(B_1 \wedge B_2)$ for simple $B_1, B_2$, distinguishing single-plane simple
    rotations from multi-plane double rotations in $n \ge 4$.
  \item \textbf{Grade 0 ($-\langle B_1, B_2 \rangle_{\Euc}$):} Controls Cauchy--Schwarz equality
    $\|B_1 + B_2\|_{\Euc} = \|B_1\|_{\Euc} + \|B_2\|_{\Euc}$, which requires positive collinearity
    $u_1 = u_2$.
\end{itemize}

\section{Positive Theorem in Dimension \texorpdfstring{$n = 2$}{n = 2} and Coplanar Regimes}

When $n = 2$ (or whenever rotations are restricted to a single oriented $2$-plane $u$ inside
$\mathbb{R}^n$), $\dim \bigwedge^2 \mathbb{R}^2 = 1$ and $\bigwedge^4 \mathbb{R}^2 = \{0\}$,
collapsing both higher-dimensional obstructions and validating every clause of the conjecture.

\begin{theorem}[Planar Abelian Validity and Non-Vacuity]\label{thm:n2}
Let $R$ be any commutative ring and let $\Cl_{2,0}(R)$ be the $4$-dimensional Euclidean
Clifford algebra with basis $\{1, e_1, e_2, e_{12}\}$. For all $a_1, a_2, c_1, s_1, c_2, s_2 \in R$:
\begin{enumerate}
  \item \textbf{Universal Lie bracket identity:} $[a_1 e_{12}, a_2 e_{12}] = 0$ holds identically.
  \item \textbf{Exact rotor angle-addition law:} Setting $\mathcal{R}(c,s) = c - s e_{12}$,
    \begin{equation}
      \mathcal{R}(c_1, s_1)\,\mathcal{R}(c_2, s_2)
      = \mathcal{R}(c_1 c_2 - s_1 s_2,\; s_1 c_2 + c_1 s_2).
    \end{equation}
    Over $\mathbb{R}$, for $B_k = -\frac{\theta_k}{2} e_{12}$, $\exp(B_1)\exp(B_2) =
    \exp(B_1 + B_2) = R(\theta_1 + \theta_2, e_{12})$.
  \item \textbf{Spectral angle additivity:} For $B_1 = -a_1 e_{12}$ and $B_2 = -a_2 e_{12}$,
    $(B_1 + B_2)^2 = -(a_1 + a_2)^2 \cdot 1$, and on $R^2$ the adjoint operator satisfies
    $\ad_{B_1 + B_2}^2(v) = -4(a_1 + a_2)^2 v$ for all $v \in R^2$.
  \item \textbf{Rational non-vacuity:} In $\Spin(2, \mathbb{Q})$, the non-central Pythagorean
    rotors $R_1 = \frac{4}{5} - \frac{3}{5}e_{12}$ and $R_2 = \frac{12}{13} - \frac{5}{13}e_{12}$
    satisfy $R_k \rev{R_k} = 1$, $R_k \neq \pm 1$, and $R_1 R_2 = \frac{33}{65} - \frac{56}{65}e_{12}$.
\end{enumerate}
\end{theorem}

\begin{proof}
Since $e_{12}^2 = e_1 e_2 e_1 e_2 = -e_1^2 e_2^2 = -1$, the even subalgebra
$\Cl_{2,0}^+(R) = R \cdot 1 \oplus R \cdot e_{12}$ is isomorphic to $R[i]/(i^2 + 1)$ and is
commutative, yielding $[a_1 e_{12}, a_2 e_{12}] = 0$. Expanding
$(c_1 - s_1 e_{12})(c_2 - s_2 e_{12}) = c_1 c_2 - (s_1 c_2 + c_1 s_2)e_{12} + s_1 s_2 e_{12}^2$
with $e_{12}^2 = -1$ gives (2); over $\mathbb{R}$, substituting $c_k = \cos(\theta_k/2)$ and
$s_k = \sin(\theta_k/2)$ applies the trigonometric addition formulas. Next,
$(B_1 + B_2)^2 = (-(a_1 + a_2)e_{12})^2 = -(a_1 + a_2)^2$. For $v = x e_1 + y e_2$,
$e_{12} e_1 - e_1 e_{12} = -2 e_2$ and $e_{12} e_2 - e_2 e_{12} = 2 e_1$, so
$\ad_{-(a_1+a_2)e_{12}}(x e_1 + y e_2) = -2(a_1 + a_2) y e_1 + 2(a_1 + a_2) x e_2$, whose square
is $-4(a_1 + a_2)^2 (x e_1 + y e_2)$. Finally, $(4/5)^2 + (3/5)^2 = 1$,
$(12/13)^2 + (5/13)^2 = 1$, and $\frac{4\cdot 12 - 3\cdot 5}{65} = \frac{33}{65}$,
$\frac{3\cdot 12 + 4\cdot 5}{65} = \frac{56}{65}$ establish (4).
\end{proof}

\section{Master Identities and Exact Characterizations in \texorpdfstring{$\Spin(3)$}{Spin(3)}}

In dimension $n = 3$, $\bigwedge^4 \mathbb{R}^3 = \{0\}$ so every bivector is simple, but
$\dim \spinlie(3) = 3 > 1$ with orthonormal basis $I = e_{23}$, $J = e_{31}$, $K = e_{12}$
satisfying $I^2 = J^2 = K^2 = -1$ and $I J = -K$, $J K = -I$, $K I = -J$.

\begin{theorem}[Master Sum-of-Squares Laws and Exact Characterizations in $\Cl_{3,0}$]\label{thm:spin3-iff}
Let $R$ be a strict ordered commutative ring (such as $\mathbb{Z}$, $\mathbb{Q}$, or
$\mathbb{R}$), and let $u_1, u_2 \in \spinlie(3, R)$ be unit bivectors
($\|u_1\|_{\Euc}^2 = \|u_2\|_{\Euc}^2 = 1$).
\begin{enumerate}
  \item \textbf{Lagrange identity and commutativity:} For all $B_1, B_2 \in \spinlie(3, R)$,
    \begin{equation}\label{eq:lagrange}
      \|[B_1, B_2]\|_{\Euc}^2
      = 4\bigl(\|B_1\|_{\Euc}^2 \|B_2\|_{\Euc}^2 - \langle B_1, B_2 \rangle_{\Euc}^2\bigr).
    \end{equation}
    In particular, $[e_{12}, e_{23}] = 2 e_{13} = -2 e_{31} \neq 0$, whereas the Jacobi identity
    $[B_1, [B_2, B_3]] + [B_2, [B_3, B_1]] + [B_3, [B_1, B_2]] = 0$ holds universally. For unit
    bivectors, $[u_1, u_2] = 0$ \textbf{if and only if} $u_1 = u_2$ or $u_1 = -u_2$.
  \item \textbf{Rotor composition master law:} For $R_k = c_k - s_k u_k$ ($c_k, s_k \in R$),
    \begin{align}
      2\bigl(\grade{R_1 R_2}{0} - (c_1 c_2 - s_1 s_2)\bigr)
      &= s_1 s_2 \|u_1 - u_2\|_{\Euc}^2, \label{eq:master-scalar}\\
      2\grade{R_1 R_2}{2}
      &= -2 s_1 c_2 u_1 - 2 c_1 s_2 u_2 + s_1 s_2 [u_1, u_2]. \label{eq:master-grade2}
    \end{align}
    Whenever $s_1 \neq 0$ and $s_2 \neq 0$, $\grade{R_1 R_2}{0} = c_1 c_2 - s_1 s_2$ holds
    \textbf{if and only if} $u_1 = u_2$. Moreover, whenever $s_1 \neq 0$, $s_2 \neq 0$, and
    $[u_1, u_2] \neq 0$, the composite rotor satisfies
    \begin{equation}\label{eq:subalg-obstruction}
      R_1 R_2 \notin R \cdot 1 \oplus \mathrm{span}_R(u_1, u_2).
    \end{equation}
    Over $\mathbb{R}$, since $(B_1 + B_2)^2 = -\|B_1 + B_2\|_{\Euc}^2 \cdot 1 \in \mathbb{R}$,
    $\exp(B_1 + B_2) \in \mathbb{R}\cdot 1 \oplus \mathrm{span}_{\mathbb{R}}(u_1, u_2)$; thus
    \eqref{eq:subalg-obstruction} implies $\exp(B_1)\exp(B_2) \neq \exp(B_1 + B_2)$ whenever
    $\sin(\theta_1/2) \neq 0$, $\sin(\theta_2/2) \neq 0$, and $u_1 \neq \pm u_2$.
  \item \textbf{Spectral sum master law:} For $B_k = -a_k u_k$ ($a_1, a_2 \in R$),
    $(B_1 + B_2)^2 = -\|B_1 + B_2\|_{\Euc}^2 \cdot 1$, $\ad_{B_1 + B_2}^3 +
    4\|B_1 + B_2\|_{\Euc}^2 \ad_{B_1 + B_2} = 0$ on $R^3$, and
    \begin{equation}\label{eq:master-spectral}
      (a_1 + a_2)^2 - \|B_1 + B_2\|_{\Euc}^2 = a_1 a_2 \|u_1 - u_2\|_{\Euc}^2.
    \end{equation}
    Whenever $a_1 \neq 0$ and $a_2 \neq 0$, spectral angle additivity
    $\|B_1 + B_2\|_{\Euc}^2 = (a_1 + a_2)^2$ holds \textbf{if and only if} $u_1 = u_2$.
\end{enumerate}
\end{theorem}

\begin{proof}
Writing $B_k = x_k e_{23} + y_k e_{31} + z_k e_{12}$, direct expansion gives
$[B_1, B_2] = -2(y_1 z_2 - z_1 y_2)e_{23} - 2(z_1 x_2 - x_1 z_2)e_{31}
- 2(x_1 y_2 - y_1 x_2)e_{12}$, so \eqref{eq:lagrange} is Lagrange's polynomial identity. Over a
strict ordered commutative ring, a sum of squares vanishes if and only if each summand vanishes.
For unit bivectors $\|u_1\|_{\Euc}^2 = \|u_2\|_{\Euc}^2 = 1$, we have
$\|u_1 \mp u_2\|_{\Euc}^2 = 2(1 \mp \langle u_1, u_2 \rangle_{\Euc})$. Hence $[u_1, u_2] = 0
\iff \langle u_1, u_2 \rangle_{\Euc}^2 = 1 \iff \langle u_1, u_2 \rangle_{\Euc} = \pm 1 \iff
u_1 = \pm u_2$.

Next, by Lemma~\ref{lem:grade-decomp} with $u_1 \wedge u_2 = 0$,
$R_1 R_2 = (c_1 - s_1 u_1)(c_2 - s_2 u_2) = (c_1 c_2 - s_1 s_2 \langle u_1, u_2 \rangle_{\Euc})
- s_1 c_2 u_1 - c_1 s_2 u_2 + \frac{1}{2} s_1 s_2 [u_1, u_2]$.
Subtracting $c_1 c_2 - s_1 s_2$ from $\grade{R_1 R_2}{0}$ gives
$s_1 s_2 (1 - \langle u_1, u_2 \rangle_{\Euc}) = \frac{1}{2} s_1 s_2 \|u_1 - u_2\|_{\Euc}^2$,
which is \eqref{eq:master-scalar}. Since $R$ is an integral domain of characteristic $0$,
for $s_1, s_2 \neq 0$ the right-hand side vanishes if and only if $\|u_1 - u_2\|_{\Euc}^2 = 0$,
i.e.\ $u_1 = u_2$. Furthermore, $\langle u_k, [u_1, u_2] \rangle_{\Euc} = 0$ for $k = 1, 2$ by
determinant antisymmetry, so every $Y = \alpha + \beta u_1 + \gamma u_2$ satisfies
$\langle \grade{Y}{2}, [u_1, u_2] \rangle_{\Euc} = 0$, whereas
$2\langle \grade{R_1 R_2}{2}, [u_1, u_2] \rangle_{\Euc} = s_1 s_2 \|[u_1, u_2]\|_{\Euc}^2 \neq 0$
whenever $s_1, s_2 \neq 0$ and $[u_1, u_2] \neq 0$, proving \eqref{eq:subalg-obstruction}.

Finally, $\|B_1 + B_2\|_{\Euc}^2 = a_1^2 + a_2^2 + 2 a_1 a_2 \langle u_1, u_2 \rangle_{\Euc}$,
so $(a_1 + a_2)^2 - \|B_1 + B_2\|_{\Euc}^2 = 2 a_1 a_2(1 - \langle u_1, u_2 \rangle_{\Euc})
= a_1 a_2 \|u_1 - u_2\|_{\Euc}^2$, establishing \eqref{eq:master-spectral}, while expanding
$\ad_B^3(v) + 4\|B\|_{\Euc}^2 \ad_B(v)$ componentwise yields zero identically.
\end{proof}

\begin{remark}[Strict Spherical Triangle Inequality]
Over $\mathbb{R}$, when $\theta_1, \theta_2 \in (0, \pi]$ and $u_1 \neq u_2$,
$\sin(\theta_1/2)\sin(\theta_2/2) > 0$ and $\|u_1 - u_2\|_{\Euc}^2 > 0$, so
\eqref{eq:master-scalar} gives $\cos(\Theta/2) > \cos((\theta_1 + \theta_2)/2)$, hence
$\Theta < \theta_1 + \theta_2$, and \eqref{eq:master-spectral} gives
$2\|B_1 + B_2\|_{\Euc} < \theta_1 + \theta_2$. Even in the anti-aligned coplanar case
$u_2 = -u_1$ (where $[B_1, B_2] = 0$ holds), $\|u_1 - u_2\|_{\Euc}^2 = 4 \neq 0$ and the
unsigned composite angle is $|\theta_1 - \theta_2| < \theta_1 + \theta_2$ (which equals the
algebraic sum only when signed angles are measured relative to the single common plane $u_1$).
\end{remark}

\section{Explicit Counterexamples Across All Dimensions \texorpdfstring{$n \ge 3$}{n >= 3}}

\subsection{Orthogonal \texorpdfstring{$90^\circ$}{90-degree}, \texorpdfstring{$180^\circ$}{180-degree}, and Rational Rotors in \texorpdfstring{$\Spin(3)$}{Spin(3)}}

\begin{proposition}[Three-Way Angle Separation for Orthogonal $90^\circ$ Rotors]\label{prop:90deg}
In $\Cl_{3,0}(\mathbb{R})$ (and algebraically over $\mathbb{Q}[\sqrt{2}]$), consider the two
orthogonal $90^\circ$ ($\theta_1 = \theta_2 = \pi/2$) unit rotors
\begin{equation}
  R_1 = \exp\!\left(-\frac{\pi}{4} e_{12}\right) = \frac{1 - e_{12}}{\sqrt{2}}, \qquad
  R_2 = \exp\!\left(-\frac{\pi}{4} e_{23}\right) = \frac{1 - e_{23}}{\sqrt{2}},
\end{equation}
with generators $B_1 = -\frac{\pi}{4} e_{12}$ and $B_2 = -\frac{\pi}{4} e_{23}$. Then:
\begin{enumerate}
  \item \textbf{True composite rotation angle $\Theta = \frac{2\pi}{3} = 120^\circ$:}
    \begin{equation}
      R_1 R_2 = \frac{1 - e_{12} - e_{23} + e_{13}}{2}
      = \cos\frac{\pi}{3} - \sin\frac{\pi}{3}\,
        \frac{e_{12} + e_{23} - e_{13}}{\sqrt{3}}.
    \end{equation}
    Moreover, $(R_1 R_2)^3 = -1$, $(R_1 R_2)^2 \neq -1$, and $\rho(R_1 R_2)$ cyclically
    permutes $(e_1, e_2, e_3) \mapsto (e_2, e_3, e_1)$ around the axis $e_1 + e_2 + e_3$.
  \item \textbf{Generator-sum spectral angle $2\|B_1 + B_2\|_{\Euc} = \frac{\pi}{\sqrt{2}}
    \approx 127.28^\circ$:} Since $(e_{12} + e_{23})^2 = -2$, $\|B_1 + B_2\|_{\Euc} =
    \frac{\pi}{2\sqrt{2}}$, so $\Spec(\ad_{B_1 + B_2}) = \{0, \pm i\pi/\sqrt{2}\}$ and
    $\exp(B_1 + B_2) \neq R_1 R_2$.
  \item \textbf{Failure of all three clauses:} Comparing normalized squared angles in $\mathbb{Q}$,
    \begin{equation}
      \left(\frac{\Theta}{\pi}\right)^2 = \frac{4}{9} \;\neq\;
      \left(\frac{2\|B_1 + B_2\|_{\Euc}}{\pi}\right)^2 = \frac{1}{2} \;\neq\;
      \left(\frac{\theta_1 + \theta_2}{\pi}\right)^2 = 1,
    \end{equation}
    while $[B_1, B_2] = \frac{\pi^2}{8} e_{13} \neq 0$.
\end{enumerate}
\end{proposition}

\begin{proof}
Using $e_{12} e_{23} = e_1 e_2^2 e_3 = e_{13} = -e_{31}$ gives the formula for $R_1 R_2$.
Squaring yields $(R_1 R_2)^2 = \frac{1}{4}(1 - 3 - 2(e_{12} + e_{23} - e_{13}))
= \frac{-1 - e_{12} - e_{23} + e_{13}}{2} \neq -1$, and multiplying once more by $R_1 R_2$
gives $(R_1 R_2)^3 = -1$. Direct evaluation of $(R_1 R_2) e_k \rev{(R_1 R_2)}$ yields
$e_1 \mapsto e_2 \mapsto e_3 \mapsto e_1$. The remaining identities follow immediately from
$(e_{12} + e_{23})^2 = e_{12}^2 + e_{23}^2 + \{e_{12}, e_{23}\} = -2$ and
$[e_{12}, e_{23}] = 2e_{13}$.
\end{proof}

\begin{proposition}[Pure Integer and Rational Counterexamples in $\Cl_{3,0}(\mathbb{Z})$ and $\Cl_{3,0}(\mathbb{Q})$]\label{prop:rational-3d}
\leavevmode
\begin{enumerate}
  \item \textbf{$180^\circ$ orthogonal rotors over $\mathbb{Z}$:} $R_1 = -e_{12}$
    ($\theta_1 = \pi$) and $R_2 = -e_{23}$ ($\theta_2 = \pi$) satisfy
    $R_1 R_2 = e_{13} \neq \pm 1$, a rotation of angle $\Theta = \pi \neq 2\pi = \theta_1 + \theta_2$.
  \item \textbf{Pythagorean rational spectrum over $\mathbb{Q}$:} $B_1 = -\frac{3}{2}e_{12}$
    ($\theta_1 = 3$) and $B_2 = -2e_{23}$ ($\theta_2 = 4$) satisfy
    $\|B_1 + B_2\|_{\Euc}^2 = \frac{25}{4} = (\frac{5}{2})^2 \neq \frac{49}{4} = (\frac{3+4}{2})^2$,
    so $\ad_{B_1 + B_2}$ has characteristic polynomial $\lambda(\lambda^2 + 25) \neq
    \lambda(\lambda^2 + 49)$.
  \item \textbf{Cayley rational unit rotors over $\mathbb{Q}$:} $R_1 = \frac{4}{5} - \frac{3}{5}e_{12}$
    and $R_2 = \frac{12}{13} - \frac{5}{13}e_{23}$ satisfy
    $\grade{R_1 R_2}{0} = \frac{48}{65} \neq \frac{33}{65} = c_1 c_2 - s_1 s_2$.
\end{enumerate}
\end{proposition}

\begin{proof}
Immediate by exact arithmetic in $\Cl_{3,0}(\mathbb{Z})$ and $\Cl_{3,0}(\mathbb{Q})$.
\end{proof}

\subsection{The Four-Dimensional Commuting Non-Simplicity Obstruction}

One might ask whether Conjecture \texttt{00000007607} can be salvaged simply by adding the
hypothesis $[B_1, B_2] = 0$. In dimension $n \ge 4$, even that weaker statement is false!

\begin{theorem}[Commuting Disjoint-Plane Counterexample in $\Spin(4)$]\label{thm:spin4}
In $\Cl_{4,0}^+(\mathbb{R})$ (and over any strict ordered commutative ring $R$), let
$B_1 = -a_1 e_{12}$ and $B_2 = -a_2 e_{34}$ with $a_1, a_2 \neq 0$ ($\theta_k = 2a_k$). Then:
\begin{enumerate}
  \item \textbf{Vanishing Lie bracket and valid exponential law:} $[e_{12}, e_{34}] = 0$, so
    $[B_1, B_2] = 0$ and (over $\mathbb{R}$) $\exp(B_1)\exp(B_2) = \exp(B_1 + B_2)$.
  \item \textbf{Non-scalar square and chiral spinor splitting:}
    \begin{equation}
      (B_1 + B_2)^2 = -(a_1^2 + a_2^2)\cdot 1 + 2 a_1 a_2\, e_{1234}
      \neq -(a_1 + a_2)^2 \cdot 1,
    \end{equation}
    and $\bigl((B_1 + B_2)^2 + (a_1 + a_2)^2\bigr)\bigl((B_1 + B_2)^2 + (a_1 - a_2)^2\bigr) = 0$,
    splitting $\Spec(L_{B_1 + B_2})$ into $\{\pm i(a_1 + a_2),\, \pm i(a_1 - a_2)\}$.
  \item \textbf{Failure of vector-spectrum angle additivity:} On $R^4$, $A = \ad_{B_1 + B_2} \in
    \solie(4, R)$ generates an orthogonal double rotation satisfying
    $(A^2 + \theta_1^2 I)(A^2 + \theta_2^2 I) = 0$ with $\Spec(A) = \{\pm i\theta_1, \pm i\theta_2\}$.
    For $(\theta_1, \theta_2) = (3, 4)$, $A^2(A^2 + (3 + 4)^2 I)(e_1) = -360\, e_1 \neq 0$.
\end{enumerate}
\end{theorem}

\begin{proof}
Since $\{1, 2\} \cap \{3, 4\} = \emptyset$, $e_{12} e_{34} = e_1 e_2 e_3 e_4 = e_{1234}$ and
$e_{34} e_{12} = (-1)^4 e_1 e_2 e_3 e_4 = e_{1234}$, giving $[e_{12}, e_{34}] = 0$ and
$(B_1 + B_2)^2 = a_1^2 e_{12}^2 + a_2^2 e_{34}^2 + 2a_1 a_2 e_{1234} = -(a_1^2 + a_2^2) +
2a_1 a_2 e_{1234}$. Because $2 a_1 a_2 \neq 0$, the grade-$4$ component is nonzero. Using
$e_{1234}^2 = 1$, $\bigl(2a_1 a_2(1 + e_{1234})\bigr)\bigl(-2a_1 a_2(1 - e_{1234})\bigr)
= -4a_1^2 a_2^2(1 - e_{1234}^2) = 0$ gives (2). On $R^4 = \mathrm{span}(e_1, e_2) \oplus
\mathrm{span}(e_3, e_4)$, $A = \mathrm{diag}(\theta_1 J_2, \theta_2 J_2)$ with $J_2^2 = -I_2$,
so $A^2(e_1) = -\theta_1^2 e_1$; for $(\theta_1, \theta_2) = (3, 4)$,
$A^2(A^2 + 49 I)(e_1) = (-9)(-9 + 49)e_1 = -360 e_1 \neq 0$.
\end{proof}

\begin{table}[ht]
\centering
\small
\setlength{\tabcolsep}{5pt}
\renewcommand{\arraystretch}{1.25}
\begin{tabular}{|l|c|c|c|}
\hline
\textbf{Regime / Example} & \textbf{Bracket $[B_1, B_2]$} &
\textbf{Composite $\Theta$ vs.\ $\theta_1 + \theta_2$} &
\textbf{Vector $\Spec(\ad_{B_1 + B_2})$} \\
\hline
$n = 2$ ($u_1 = u_2$) & $0$ (identity) &
$\Theta = \theta_1 + \theta_2$ (exact) & $\{\pm i(\theta_1 + \theta_2)\}$ \\
\hline
$n = 3$, $90^\circ$ ($e_{12}, e_{23}$) & $\frac{\pi^2}{8} e_{13} \neq 0$ &
$\frac{2\pi}{3} \neq \pi$ ($2\|B_1+B_2\| = \frac{\pi}{\sqrt{2}}$) &
$\{0, \pm i\pi/\sqrt{2}\}$ \\
\hline
$n = 3$, $(3,4)$ ($e_{12}, e_{23}$) & $6 e_{13} \neq 0$ &
$\Theta < 7$ ($2\|B_1+B_2\| = 5$) & $\{0, \pm 5i\} \neq \{0, \pm 7i\}$ \\
\hline
$n = 4$, $(3,4)$ ($e_{12}, e_{34}$) & $0$ (commuting!) &
Double rotation $(3, 4)$ & $\{\pm 3i, \pm 4i\} \neq \{0,0,\pm 7i\}$ \\
\hline
\end{tabular}
\caption{Comparison across dimensions $n = 2, 3, 4$ showing how all three clauses of
Conjecture \texttt{00000007607} succeed for $n = 2$ ($u_1 = u_2$) and fail for $n \ge 3$.}
\label{tab:comparison}
\end{table}

\subsection{Universal Extension to Every Dimension \texorpdfstring{$n \ge 3$}{n >= 3}}

\begin{theorem}[Universal Failure for All $n \ge 3$]\label{thm:universal-n}
For every integer $n \ge 3$, the canonical isometric embedding $\mathbb{R}^3 \hookrightarrow
\mathbb{R}^n$ induces injective homomorphisms $\Cl_{3,0}(\mathbb{R}) \hookrightarrow
\Cl_{n,0}(\mathbb{R})$ and $\Spin(3) \hookrightarrow \Spin(n)$. Specifically, let
$Q_{12}^{(n)}, Q_{23}^{(n)} \in \SO(n, \mathbb{Z})$ be the $90^\circ$ rotations in the
$(e_1, e_2)$ and $(e_2, e_3)$ coordinate planes (fixing $e_4, \dots, e_n$), with unit-angle
generators $A_{12}^{(n)}, A_{23}^{(n)} \in \solie(n, \mathbb{Z})$. Then for every $n \ge 3$:
\begin{enumerate}
  \item $Q_{12}^{(n)} Q_{23}^{(n)}$ permutes $(e_1, e_2, e_3) \mapsto (e_2, e_3, e_1)$ of
    order $3$ and has trace $\tr(Q_{12}^{(n)} Q_{23}^{(n)}) = n - 3$, whereas any single-plane
    rotation of additive angle $\frac{\pi}{2} + \frac{\pi}{2} = \pi$ has trace $n - 4 \neq n - 3$.
  \item $S^{(n)} = A_{12}^{(n)} + A_{23}^{(n)}$ satisfies $S^{(n)}\bigl((S^{(n)})^2 + 2I_n\bigr) = 0$
    while $\bigl((S^{(n)})^2 + 4I_n\bigr)(e_1) \neq 0$, refuting spectral angle additivity
    $(1 + 1)^2 = 4$.
  \item $[A_{12}^{(n)}, A_{23}^{(n)}](e_1) \neq 0$, refuting universal Lie bracket vanishing.
\end{enumerate}
\end{theorem}

\begin{proof}
On $\mathrm{span}(e_1, e_2, e_3)$ the operators coincide with the $n = 3$ operators of
Proposition~\ref{prop:90deg} (scaled to unit generators), while on $\mathrm{span}(e_4, \dots, e_n)$
$Q_{12}^{(n)}$ and $Q_{23}^{(n)}$ act as $I_{n-3}$ and $A_{12}^{(n)}, A_{23}^{(n)}$ act as $0$.
Thus $\tr(Q_{12}^{(n)} Q_{23}^{(n)}) = 0 + (n - 3) = n - 3 \neq (n - 2) + 2\cos\pi = n - 4$,
$\bigl((S^{(n)})^2 + 4I_n\bigr)(e_1) = 3e_1 + e_3 \neq 0$, and
$[A_{12}^{(n)}, A_{23}^{(n)}](e_1) = -e_3 \neq 0$.
\end{proof}

\section{Formal Verification in Core Lean 4.19.0 and Python}

Every algebraic foundation, positive $n = 2$ theorem, master sum-of-squares characterization,
and $n \ge 3$ counterexample in this paper has been formally verified in Core Lean~4.19.0
(inside \texttt{lean/}, with zero external dependencies, \texttt{-DwarningAsError=true}, and
zero \texttt{sorry}, \texttt{admit}, \texttt{native\_decide}, or custom axioms), alongside an
independent exact bitmask Clifford algebra verification suite (\texttt{verify.py}).
Table~\ref{tab:lean-map} summarizes the correspondence between paper results and Lean~4
declarations.

\begin{table}[ht]
\centering
\footnotesize
\setlength{\tabcolsep}{5pt}
\renewcommand{\arraystretch}{1.22}
\begin{tabular}{|>{\raggedright\arraybackslash}p{1.62in}|l|>{\raggedright\arraybackslash}p{2.58in}|}
\hline
\textbf{Paper Result} & \textbf{Lean 4 Module} & \textbf{Verified Declaration(s)} \\
\hline
Reflective ring normalizer & \texttt{RingBase.lean} &
\texttt{CommRing.eq\_of\_toPoly\_eq}, \newline \texttt{comm\_ring} \\
\hline
Ordered rings $\mathbb{Z}, \mathbb{Q}, \mathbb{Q}[\sqrt{2}]$ & \texttt{Ring.lean} &
\texttt{StrictOrderedCommRing}, \texttt{QRat}, \texttt{QSqrt2} \\
\hline
Theorem~\ref{thm:n2} ($n = 2$ positive) & \texttt{Clifford2.lean} &
\texttt{rotor2\_mul\_angle\_add}, \newline \texttt{spin2\_positive\_nonvacuity} \\
\hline
Theorem~\ref{thm:spin3-iff}(1) (Lagrange / bracket) & \texttt{Clifford3.lean} &
\texttt{bivec\_bracket\_normSq\_lagrange}, \newline \texttt{unit\_bivec\_commute\_iff} \\
\hline
Theorem~\ref{thm:spin3-iff}(2) (Rotor master / span) & \texttt{Clifford3.lean} &
\texttt{rotor\_scalar\_angle\_add\_iff}, \newline \texttt{rotor\_mul\_not\_in\_gen\_span} \\
\hline
Theorem~\ref{thm:spin3-iff}(3) (Spectral / adjoint) & \texttt{Clifford3.lean} &
\texttt{gen\_sum\_spectral\_angle\_add\_iff}, \newline \texttt{adjoint\_cayley\_hamilton} \\
\hline
Propositions~\ref{prop:90deg}--\ref{prop:rational-3d} & \texttt{Clifford3.lean} &
\texttt{counterexample\_90deg\_orthogonal}, \newline \texttt{counterexample\_pythagorean\_rat} \\
\hline
Theorem~\ref{thm:spin4} ($\Spin(4)$ commuting) & \texttt{HigherDim.lean} &
\texttt{spin4\_commuting\_still\_not\_angle\_additive} \\
\hline
Theorem~\ref{thm:universal-n} (All $n \ge 3$) & \texttt{HigherDim.lean} &
\texttt{DimN.universal\_refutation\_ge\_3} \\
\hline
Top-level synthesis & \texttt{Conjecture00000007607.lean} &
\texttt{positive\_planar\_spin2}, \newline \texttt{conjecture\_00000007607\_disproved} \\
\hline
\end{tabular}
\caption{Mapping between mathematical results and machine-checked Core Lean~4.19.0 theorems.}
\label{tab:lean-map}
\end{table}

\begin{thebibliography}{9}
\bibitem{DoranLasenby}
C.~Doran and A.~Lasenby, \emph{Geometric Algebra for Physicists}, Cambridge University Press,
2003.

\bibitem{HestenesSobczyk}
D.~Hestenes and G.~Sobczyk, \emph{Clifford Algebra to Geometric Calculus: A Unified Language
for Mathematics and Physics}, D.~Reidel Publishing, 1984.

\bibitem{Lounesto}
P.~Lounesto, \emph{Clifford Algebras and Spinors}, 2nd ed., London Mathematical Society Lecture
Note Series, vol.~286, Cambridge University Press, 2001.

\bibitem{Hall}
B.~C.~Hall, \emph{Lie Groups, Lie Algebras, and Representations: An Elementary Introduction},
2nd ed., Graduate Texts in Mathematics, vol.~222, Springer, 2015.
\end{thebibliography}

\end{document}
